\chapter[A World More Alike or Different?]{Is the World Becoming More Alike or More Different?}
\section{A Football's Birthplace}
Pick up an ordinary football, its thirty-two panels stitched into a black-and-white pattern.

Read its markings. The famous sports brand has headquarters in America or Germany. FIFA certification comes from an international organization based in Switzerland. The star your classmates idolize is Argentine, now playing for an American club, while the platform broadcasting him belongs to someone from another country.

Turn the ball over and find its origin in small print. A good hand-stitched ball may well come from somewhere unfamiliar: Sialkot, an industrial city in northeastern Pakistan. Most hand-stitched footballs originate in this city and surrounding towns. World Cup balls, European clubs' training balls, and your school's peeling old ball may pass through the same workers' hands. Under lamps, they align panels, stitch inside out, and close seams. Their piece-rate wages and the eventual retail price stand a whole world apart.

The same ball is your toy, Sialkot's job, a brand's spreadsheet entry, and an Argentine slum child's imagined escape. Like its thirty-two panels, it stitches you to distant strangers.

Goods, people, money, videos, and music move faster worldwide than ever. Does this make the world more alike---same brands, songs, and lives---or more different? Or do both claims mislead?

\section{Eight at Night in Yiwu}
Enter a scene with me.

It is eight in Yiwu, Zhejiang, a city renowned for foreign trade. Its International Trade Mart, the world's largest small-commodities market, joins tens of thousands of stalls. Day belongs to business; now belongs to dinner.

On a nearby street, Arabic, English, and Russian signs mingle with Chinese. Grilled-meat smoke and cardamom escape a Yemeni restaurant. Four men near the window argue in Arabic about a shipment's arrival over tea refilled three times. Next door, a Somali trader video-calls family in Mogadishu in Swahili, children's laughter reaching his raised phone. At the entrance, Chinese staff clear tables and promise departing guests a discount in accented English.

Passersby come from more than a dozen places: Yemenis, Iraqis, and Syrians brought by war and business; Nigerian and Senegalese buyers hauling sample cases; Indian interpreters switching hundreds of times daily among Hindi, English, and Chinese; traders from across China; local people whose parents often began peddling candy for chicken feathers. An Egyptian regular once told a visitor, in paraphrase, that books in Cairo introduced China, but a Yiwu dinner table made it real.

Notice three components of this scene.

First, goods flow like water: hair ties leave Yiwu and appear three months later in Lagos, Nigeria's largest city, at half local prices. Second, people move less easily: the Yemeni owner's visa, residence, and children's schooling create greater difficulties than containers. Third, meaning travels and changes. Yiwu-made Santas hang on Lima Christmas trees, while the street's favorite Yemeni grill quietly lowers its heat for Chinese customers.

Goods, people, and meanings cross borders under different speeds, routes, and rules. Remember the difference.

\section{Is the World Photocopying Itself?}
Enlarge Yiwu's miniature world, and a century-old dispute appears.

Evidence of sameness lies everywhere.

Young people in Moscow and Mexico City wear similar hoodies, jeans, and sneakers. One song charts in dozens of countries in a week; classmates hum what peers in Lagos and Lima hum. Airports resemble siblings, with identical duty-free stores, coffee chains, and sign colors. New urban districts likewise repeat glass facades, broad roads, and chain hotels. Modern travel's disappointment, the joke goes, is flying for hours and landing as though you never left.

Migration and tourism move unprecedented numbers. Hundreds of millions live outside their birth countries, roughly one person in thirty. International tourist arrivals rose from over twenty million annually in the mid-twentieth century to over a billion before the pandemic. Seeing elsewhere became mass consumption; being seen became business. Towns redecorate according to tourists' imagined authenticity, becoming what visitors expect. Revenue arrives alongside the risk of culture turning into scenery.

Evidence of difference is equally plentiful.

Indian audiences want song-and-dance additions to Hollywood films. One fried-chicken brand sells soy milk and fried dough in China, vegetarian burgers in India. Fans share a World Cup but Senegalese victory celebrations and Japanese post-defeat gestures inhabit different worlds. One powerful recent musical wave comes from Korea, whose idol industry succeeded by becoming \textbf{different} from American pop. Against everyone-drinks-cola claims, Peru's yellow Inca Kola became a national devotion. Coca-Cola could not beat it locally for years and eventually bought it.

The real question is therefore not a binary choice:

\textbf{Which parts become more alike, which more different, and who decides? When cultures meet, do they copy, mix, or swallow one another?}

Behind this stands a sharper question: whom does resemblance benefit? Do Sialkot stitchers, Yiwu restaurateurs, Lagos hair-tie vendors, and you occupy equivalent places in this changing world?

Before answering, predict whether everyday life worldwide will be more similar or different fifty years hence. Keep that prediction for the end.

\section{Worries and Celebrations}
Hear concern at its strongest, not as a caricature of stubborn traditionalism.

Concern has at least three arguments. First, choice disappears. Local foods, crafts, and dialects embody generations of practical experimentation. Capital and standardization let chains displace them, leaving five worldwide options instead of more choice. Second, power is unequal. Supposedly free cultural flows emerge from a few film industries and platforms; international standards follow market share. France promoted cultural exception in trade talks, arguing films and music are not ordinary goods and deserve domestic subsidy lest Hollywood monopolize screens. Third, meaning empties out. An old town remains physically old, but skewers, Yiwu souvenirs, and chain guesthouses make it infinitely reproducible scenery. The fear is not change but everywhere becoming a tourist attraction.

These arguments are strong. You may have visited old streets from north to south seemingly stocked by one wholesaler.

Now hear celebration just as fully.

First, ordinary people gain tools and stages. Nigeria's Nollywood began with cheap equipment and discs, not Hollywood wealth, and now ranks among the largest producers by film count, telling Nigerian stories across Africa and its diaspora. Global technology carries intensely local content. Second, encounters create new, sometimes richer things: Japanese-Peruvian chefs combine Japanese techniques and Peruvian ingredients into internationally celebrated Nikkei cuisine, impossible without encounter. Third, for marginalized people, resembling the world can liberate. A small-town girl discovering many possible lives through videos gains options, not loses them.

Now hear mobile people themselves, often belonging to neither camp.

On Sundays, Hong Kong's Central walkways and squares fill with women picnicking, singing, and video-calling: Filipino domestic workers. They live with employers and care for others' children while their own remain with grandmothers in the Philippines. Many serve one household for over a decade, switching English, Cantonese, and Tagalog and remitting money monthly. Ask whether the world shrank and they may answer practically: small enough to fly thousands of kilometers for work, but still divided. Their wages, visas, and children's childhoods have entirely different dimensions from employers'.

Dubai takes the global city to an extreme: foreigners form the vast majority, citizens a minority. South Asian construction workers, globally recruited office professionals, and shoppers from over a hundred countries mingle. Borderless? Borders are unusually clear here. Passport color influences queues, wages, family reunification, and permission to remain in old age. An Indian engineer and Nepali builder may spend ten years in one city yet inhabit different Dubais.

American sociologist Ruth Hill Useem named another group third-culture kids: children growing across countries with parents, such as a Korean father's child attending international school while Mom works in Africa. Asked where they come from, many hesitate. Home lies not on one map but among airports, video calls, and similarly mobile friends. Cross-border identity is gift and burden: easy cultural switching without certainty which culture is theirs.

Together these voices show that judging sameness depends considerably on your place within the flow.

\section{Thinking Bridge: Ask Three Questions}
How can you evaluate sweeping claims of increasing sameness or diversity? Ask three useful questions whenever globalization supposedly changes a place.

\textbf{First: what is moving?}

Distinguish goods, people, and meanings such as ideas, music, fashion, and institutions. Their rules differ. Since container shipping emerged in 1956, freight costs fell repeatedly, leaving a toy's shipping a fraction of its price. People face passports, visas, and income barriers. Most European travel needed no passport before the First World War; today passport color nearly determines someone's radius of movement. Meanings move fastest and most elusively: a song circles the globe overnight but changes meaning everywhere. Always identify the moving thing.

\textbf{Second: does change affect the surface or deeper layers?}

Compare layers. Brands, clothing, music, and architectural styles converge quickly at the surface. Eating habits, family structure, and celebrations change more slowly, often digesting new things through old practices. Right and wrong, respectability, and life's purpose change across generations and often resist. Many disputes merely pit someone pointing at identical surfaces against someone pointing at different depths.

\textbf{Third: what does the same thing mean here?}

Identical objects can mean different things. Anthropologists describe cola entering religious healing and cleansing rituals in Indigenous communities of Chiapas, southern Mexico. The same bottle is your soft drink and their ritual implement. Conversely, local tradition may be imported: Sichuan and Hunan cuisine seem unimaginable without chili, yet American-origin peppers entered China only centuries ago. Ancient spicy traditions grew from post-Age-of-Exploration global tables.

Together, these questions turn broad conclusions into checkable ones: what moves, which layer changes, what it means here. Increasing sameness proves sometimes true, sometimes false, and often too crudely asked.

Practice on your shoes. Goods move through factories, brands, and freight; few people move with them, with assembly workers remaining in factories. Meaning changes from school uniform accessory to piece-rate job to collectors' speculative trade. The surface is globally identical, but what the brand means to your parents and you differs underneath. The tool works on shoes and footballs because it is portable and reusable.

\section[Imports Two Thousand Years Ago]{Reading Evidence: Imports Two Thousand Years Ago}
You may imagine global trade began with containers and aircraft. Read the Records of the Grand Historian's ``Account of Dayuan.''

After more than a decade traveling west, Han envoy Zhang Qian reached Bactria, around modern Afghanistan, and saw two surprising objects:

\begin{quote}
Qian said, ``In Bactria I saw Qiong bamboo staffs and Shu cloth. I asked where they came from. The Bactrians replied, `Our merchants bought them in Shendu.' ''
\end{quote}
In brief, he found bamboo staffs and cloth from present-day Sichuan, purchased by Bactrian traders in Shendu, the ancient name for India.

Unpack the passage. Snow mountains, deserts, and many disconnected states lay between Sichuan and Bactria, yet goods had passed hand to hand into Central Asia ahead of people, maps, or the concept of a Silk Road. Sent to open western routes, Zhang Qian discovered merchants already ahead of envoys. Goods globalize, people follow, ideas catch up: a repeatedly renewed two-thousand-year sequence.

The same-versus-different question is old too. In the 1848 Communist Manifesto, Marx and Engels made a startlingly early observation, rendered in the standard Chinese translation as follows:

\begin{quote}
By developing the world market, the bourgeoisie has made every country's production and consumption worldwide in character.
\end{quote}
They also described national one-sidedness and narrowness becoming increasingly impossible. Over 170 years ago, someone already predicted a flattened world. Yet a century and a half later, Inca Kola survives, Nollywood thrives, and Korean songs top global charts. The prediction was half right, half wrong. Why is the next section's question.

\section{One World, Three Accounts}
Offer three explanations of simultaneous convergence and differentiation. Separate facts---global circulation of brands, songs, and goods---from explanations of their routes, participants' feelings such as Peruvians calling Inca Kola ours, and judgments about whether convergence is good. Compare explanations.

\textbf{Explanation one: the strong cover the weak, or cultural homogenization.}

Flows are mainly one-way: wealthy film industries, platforms, and chains spread products, squeezing local cultures to death or into corners. English dominates not through beauty but economic power. Ownership, standard-setting, and profit provide strong evidence. The limitation is treating receivers as empty bottles. Inca Kola survives, Chiapas makes cola a ritual implement, and Nigerians create a film universe for a fraction of Hollywood costs. Homogenization explains airports, not markets.

\textbf{Explanation two: global and local mix, or hybridization.}

Sociologist Roland Robertson introduced glocalization: global goods sell by taking local root. Companies localize; communities transform imports; their mixture creates new varieties. Japanese curry took Indian spices through the British navy and became mild, sweet, national home cooking. Hawaiian pizza was invented in Canada in 1962 by a Greek immigrant, named after a canned-pineapple brand, neither Hawaiian nor Italian.

Consider salmon sushi, an easily misread counterexample. Enthusiasts often assume an ancient Japanese tradition, yet Norwegian seafood exporters strongly promoted raw salmon with rice in the 1980s and 1990s. Supposedly pure tradition emerged from recent international marketing. Before defending tradition against globalization, inspect its birth certificate: many traditions are younger than globalization.

Hybridization's limit is cheerful mixing that conceals unequal capacity to mix. Multinationals adapting taste are innovative; football stitchers seeking changed futures have no comparable brand to play.

\textbf{Explanation three: follow profit and borders, a structural account.}

Instead of cultural depth, follow power and interests. Profitable standardization produces convergent airports, supply chains, office software, and vaccine standards. Unprofitable or identity-sensitive domains retain varied marriage customs, funerals, dialects, and faiths. States loosen goods' borders while tightening people's: containers clear more easily, passports receive closer inspection. This explains both sameness and difference, including Dubai's divided queues. Its limit is coldness: treating culture as profit's shadow misses people fighting for uneconomic differences, endangered languages, or scarcely attended opera. Some human commitments defy balance sheets.

Each account sees part: homogenization sees power, hybridization creativity, structural analysis rules. Together they give a three-dimensional answer.

\section[What Dies with a Language?]{Deeper Challenge: What Dies with a Language?}
Language is the most sensitive gauge of sameness and difference.

Roughly seven thousand languages remain in use, unevenly distributed. A few dozen serve most people; thousands have only thousands, hundreds, or dozens of speakers. Papua New Guinea alone has over eight hundred, Earth's densest concentration. UNESCO estimates nearly half face disappearance, with roughly one falling silent every two weeks as its last speaker dies.

Does only vocabulary die? Northern Australia's Guugu Yimithirr offers an example. Rather than body-centered left and right, it uses cardinal directions: the cup north of you. Speakers develop remarkable orientation, reporting absolute directions even inside unfamiliar buildings. Linguists see that language does more than describe: it determines what one \textbf{must} notice. Losing a language loses a division of observational labor: plant categories, waves' meanings, ways of recounting family histories, much never written.

Distinguish two mechanisms of language death. Transmission may be forcibly broken, as Indigenous children in Canada and Australia were taken to residential schools and prohibited from speaking their languages. This is harm, not natural evolution, with claims to apology and reparation. Or families may choose usefulness, teaching children only a dominant language for school and jobs. Outsiders should not casually condemn sacrificing children's prospects to museum-like preservation. Yet calling language loss purely natural and unregrettable is too quick. Communities can teach both dominant and ancestral languages. Maori language nests, begun in 1982, immersed preschoolers and restored daily use to a shrinking language. Hebrew provides a larger precedent: long chiefly religious and written, it became street and kitchen speech again in twentieth-century Palestine, history's only revival on such a scale. Languages can die or revive, depending on communities' opportunity, resources, and dignity.

China faces its own test. Manchu, once an official Qing language leaving vast archives, now has few everyday native speakers; readers and writers are often specialists. Widely spoken Shanghai, Cantonese, and Southern Min varieties also face generational decline, with fluent elders and children who only understand. Classroom inclusion and compatibility with a common language remain disputed. There are no villains in the dilemma: a shared language connects over a billion, while local varieties hold humor and memory.

Take this challenge: language and biological diversity are usefully compared, but not identical. Species cannot choose their fate; language communities can. Ask not simply whether every language must be rescued, but \textbf{whether each community genuinely can choose}. Does a child relinquish a mother tongue freely for a wider world, or trade part of the self for admission?

\section{Today: The Global Village in Your Phone}
The old question now operates in your pocket.

One short-video dance spreads through Seoul, Sao Paulo, and your classroom, altered locally in gestures and dialect lyrics. Hybridization is algorithms' daily work. China's Mixue reaches Jakarta streets, with thousands of shops and many Indonesian young people assuming it local. Translated Chinese web novels have devoted North American readers demanding updates. Latin American K-pop fans teach themselves Korean, with Mexico City classes hard to enter. Culture no longer travels only west to east.

At the same time, people turn back toward hanfu, Chinese-inspired design, and dialect rap. Tighter global connection intensifies who-am-I questions. Such traditions are often themselves new: contemporary enthusiasts research and boldly alter historical clothing. Today's hanfu is less restoration than young people recreating tradition through global-era aesthetics. Tradition is not an antique rescued from globalization, but a play restaged upon it.

Machine translation introduces another variable. If phones instantly translate dozens of languages, why study one laboriously? The tool transmits information, but not Guugu Yimithirr's cardinal sensitivity or a language's humor, taboos, and tact. It removes information boundaries while revealing meaning's boundaries more clearly.

Recent years also teach reversal. Pandemic closures and competition for supplies dissolved borderless-world illusions overnight. Globalization proved not an unalterable fact but a condition, able to accelerate, pause, or partly reverse. Unrestricted movement of goods and people is not history's default.

\section[Which Differences Deserve Protection?]{For You: Which Differences Deserve Protection?}
Return to the football.

It stitches workers, shareholders, Yiwu traders, Filipino domestic workers, third-culture children, and you into one net. Airports, brands, dances, and vaccine standards converge, while languages, funerals, faiths, and imagined good lives stubbornly differ. Convergence is not all conspiracy, nor difference all romance.

The final question is a reasoned choice, not a prediction of resemblance:

\textbf{If some similarities deserve welcome---lifesaving vaccine standards, aviation safety, girls' schooling---and some differences deserve protection---languages, festivals, cuisines---who decides which? If a protected difference harms women or minorities, does it still deserve protection?}

Answer with reasons. Weigh three worries, three celebrations, mobile people's unequal dimensions, the three-question tool, three explanations, what dies with languages, and Maori revival.

Finally retrieve your opening prediction about fifty years hence. Now you should be able to explain it. Your reasons are themselves evidence that the world has moved and mixed once more.
